\chapter{Solutions}

\runinhead{Calcul numérique d'une
probabilité}\phantomsection\addcontentsline{toc}{subsubsection}{Calcul
numérique d'une probabilité}\label{V-answer-mc1}
cf.~notebook

\runinhead{Exemples d'application
1}\phantomsection\addcontentsline{toc}{subsubsection}{Exemples
d'application 1}\label{V-answer-exemples1}
On suppose que \(U \sim \mathcal{U}_{]0,1[}\)

\begin{itemize}
\item
  Uniforme sur un intervalle \(I \subset \mathbb{R}\)~:

  Si \(I = ]a,b[,~a< b\), alors la fonction de répartition de
  \(X \sim \mathcal{U}_I\) est
  \(F_X(x) = \frac{x-a}{b-a} 1_{]a,b[}(x) + 1_{[b,+\infty[}(x)\), d'où
  par inversion \(X = a + (b-a)U\)
\item
  Exponentielle de paramètre \(\lambda \in \mathbb{R}_+^\ast\)~:

  la fonction de répartition de \(X \sim \mathcal{E}(\lambda)\) est
  \(F_X (x) = 1-\exp(-\lambda x)\), d'où par inversion
  \(X = -\frac{\ln(1-U)}{\lambda}\). Pour simplifier, on remarquera que
  si \(U \sim \mathcal{U}_{]0,1 [ }\) , alors
  \(1- U \sim \mathcal{U}_{]0,1[}\), d'où
  \(X = -\frac{\ln(U)}{\lambda}\)
\item
  de Cauchy, de densité
  \(x\in\mathbb{R}\mapsto \left(\pi\left(1+x^2\right)\right)^{-1}\),

  la loi de Cauchy admet la fonction de répartition
  \(F(x) = \frac{1}{\pi} \arctan (x) + \frac{1}{2}\), d'où
  \(X = \tan(\pi U -1/2)\)
\item
  de Laplace de paramètres \(\mu \in \mathbb{R}\) et
  \(s\in\mathbb{R}_+^\ast\), de densité
  \(x\in\mathbb{R}\mapsto \frac{1}{2s}\,\exp\left\{-\frac{|x-\mu|}{s}\right\}\),

  la loi de Laplace admet la fonction de répartition
  \[F(x) = \left\{ \begin{array}{ll}
    \frac{1}{2} \exp(x) & \text{ si } x < 0,\\
    1 - \frac{1}{2} \exp(-x) & \text{ si } x \geq 0,
    \end{array} \right. = \frac{1}{2}\left(1 + \text{sgn}(x)(1-\exp(-|x|))\right)\]

  d'où \(X = \text{sgn}(U - 1/2)\ln(1-2|U-1/2|)\), où sgn est la
  fonction signe.
\item
  Logistique de paramètres \(\mu \in \mathbb{R}\) et
  \(s\in\mathbb{R}_+^\ast\), de fonction de répartition
  \(x\in\mathbb{R}\mapsto \left(1 + \exp\left\{-\frac{x-\mu}{s} \right\}\right)^{-1}\)~:

  par inversion,
  \(X = s \left(-\ln\left(\frac{1}{U}-1\right)+ \mu\right)\)
\end{itemize}

\runinhead{Pile ou
face}\phantomsection\addcontentsline{toc}{subsubsection}{Pile ou
face}\label{V-answer-pf}
Il suffit d'associer à ``pile'' un événement portant sur
\(U \sim \mathcal{U}_{]0,1[}\) de la bonne probabilité, \(p \in ]0,1[\),
par exemple \(\{U \leq p\}\). On obtient ainsi l'algorithme~:

\begin{enumerate}
\def\labelenumi{\arabic{enumi}.}
\tightlist
\item
  Générer \(u\) selon \(\mathcal{U}_{]0,1[}\)
\item
  Renvoyer ``pile'' si \(u \leq p\), ``face'' sinon
\end{enumerate}

\runinhead{Exemples d'application
2}\phantomsection\addcontentsline{toc}{subsubsection}{Exemples
d'application 2}\label{V-answer-exemples2}

\begin{itemize}
\item
  Binomiale de paramètres \(n\in\mathbb{N}^\ast\) et \(p \in\, ]0,1[\),

  La fonction de répartition de la loi binomiale est donnée par
  \[F(x) = \begin{cases}
    0&\text{si } \;x<0\\
    \sum_{{k=0}}^{{\lfloor x\rfloor }}{n \choose k}p^{k}(1-p)^{{n-k}} & \text{si } \;0\leq x<n\\
    1&\text{si} \;x \geq n
    \end{cases}\]

  où \({\lfloor x\rfloor }\) est la partie entière de \(x\).\\
  On peut alors découper \(]0,1[\) en intervalles de la forme
  \(]\sum_{{k=0}}^{{m}}{n \choose k}p^{k}(1-p)^{{n-k}},\sum_{{k=0}}^{{m+1}}{n \choose k}p^{k}(1-p)^{{n-k}}[\)
  pour \(m \in \{0,n-1\}\) et renvoyer \(m+1\) si un \(u\) généré selon
  \(\mathcal{U}_{] 0,1[}\) appartient à l'un de ces intervalles, 0
  sinon.

  Il est cependant beaucoup plus simple de remarquer qu'une variable
  aléatoire de loi \(\mathcal{B}(n,p)\) s'obtient comme la somme de
  \(n\) variables de Bernoulli que l'on peut simuler comme
  \protect\hyperref[V-answer-pf]{ci-dessus (p.~\pageref*{V-answer-pf})}.
\end{itemize}

\begin{itemize}
\item
  Uniforme sur l'union de deux segments non vides et disjoints
  \([a,b], [c,d]\subset\mathbb{R}\), de densité
  \(x\in\mathbb{R}\mapsto (b-a + d-c)^{-1}\,1_{[a,b]\cup[c,d]}(x)\)

  La fonction de répartition d'une telle variable est
  \[F(x) = \frac{x-a}{b-a + d-c} 1_{]a,b[}(x) + \frac{b-a}{b-a + d-c} 1_{[b,c]}(x) + \frac{x-c+b-a}{b-a + d-c} 1_{]c,d[}(x) + 1_{[d,+\infty[}(x)\]
  Pour simuler cette variable, on peut donc utiliser l'algorithme
  suivant~:

  \begin{enumerate}
  \def\labelenumi{\arabic{enumi}.}
  \tightlist
  \item
    Générer \(u\) selon \(\mathcal{U}_{]0,1[}\)
  \item
    Retourner \[\begin{cases} 
      a + u(b-a + d -c) & \text{ si } u \in ]0,\frac{b-a}{b-a + d-c}[ \\
      c + (u(b-a + d -c) - (b-a)) \text{ sinon }
      \end{cases}\]
  \end{enumerate}
\end{itemize}

\runinhead{Loi uniforme sur un
pavé}\phantomsection\addcontentsline{toc}{subsubsection}{Loi uniforme
sur un pavé}\label{V-answer-unipave}
On note \([a_1,b_1]\times\dots\times[a_d,b_d] \subset \mathbb{R}^d\). On
peut voir que la densité d'une v.a. uniforme sur un tel pavé s'écrit
\(f(x) = \frac{1}{b_1 - a_1} \cdots \frac{1}{b_d - a_d} 1_{[a_1,b_1]\times\dots\times[a_d,b_d]}(x)\)
soit comme le produit des densités de ses coordonnées qui sont donc
indépendantes. On peut donc simuler la variable d'intérêt avec cet
algorithme~:

\begin{enumerate}
\def\labelenumi{\arabic{enumi}.}
\tightlist
\item
  Générer \(u_1,\ldots,u_d\) indépendamment selon
  \(\mathcal{U}_{]0,1[}\)
\item
  Retourner \((a_1 + u_1(b_1-a_1), \ldots, a_d + u_d(b_d-a_d))\)
\end{enumerate}

\runinhead{Démonstration}\phantomsection\addcontentsline{toc}{subsubsection}{Démonstration}\label{V-answer-unicond}
Soit \(C \in \mathcal{B}(B)\) (considérer les boréliens de \(B\) suffit
à caractériser la loi conditionnelle à l'événement \(U \in B\) puisque
si \(C \cap B = \varnothing\), alors nécessairement,
\(\mathbb{P}(U \in C |U \in B) = 0\)). On a \begin{align*}
\mathbb{P}(U \in C |U \in B) & = \frac{\mathbb{P}(U\in C, U \in B)}{\mathbb{P}(U \in B)} = \frac{\mathbb{P}(U\in C)}{\mathbb{P}(U \in B)} \\
& = \frac{\int_C \frac{dx}{\lambda(A)}}{\int_B \frac{dx}{\lambda(A)}} \text{ puisque } U \sim \mathcal{U}_A \\
& = \frac{\lambda(C)}{\lambda(B)} = \int_C \frac{1}{\lambda(B)} dx
\end{align*} on reconnaît bien la loi uniforme sur \(B\).

\runinhead{Simulation}\phantomsection\addcontentsline{toc}{subsubsection}{Simulation}\label{V-answer-simunibor}
On prend un pavé \(R\) tel que \(B \subset R\) puis on applique
l'algorithme suivant~:

\begin{enumerate}
\def\labelenumi{\arabic{enumi}.}
\tightlist
\item
  générer \(u\) selon \(\mathcal{U}_R\)
\item
  Si \(u \notin R\), retour en 1.
\item
  retourner \(u\)
\end{enumerate}

\runinhead{Taux de
rejet}\phantomsection\addcontentsline{toc}{subsubsection}{Taux de
rejet}\label{V-answer-tauxrej}
Puisque le couple \((Y, Uf_Y(Y))\) est uniforme sur \(A_f\), alors le
couple \((Y, aUf_Y(Y))\) est uniforme sur \(A_{af}\), d'où
\begin{align*}
\mathbb{P}(\{aUf_Y(Y) > f_X(Y)\}) & = \mathbb{E}(1_{\{aUf_Y(Y) > f_X(Y)\}}) \\
& = \mathbb{E}(\mathbb{E}(1_{\{aUf_Y(y) > f_X(y)\}}|Y=y)) \\
& = \int_R \int_0^{af_Y(y)} 1_{\{auf_Y(Y) > f_X(Y)\}} \frac{1}{af_Y(y)} du f_Y(y) dy \\
& = \int_R \left(1 - \frac{f_X(y)}{af_Y(y)}\right) f_Y(y) dy \\
& = 1 - \frac{1}{a}
\end{align*}

\section{Loi uniforme dans un
domaine}\label{V-loi-uniforme-dans-un-domaine-1}

cf.~notebook

\section{Simulation selon la loi
géométrique}\label{V-simulation-selon-la-loi-guxe9omuxe9trique-1}

\runinhead{Question
1}\phantomsection\addcontentsline{toc}{subsubsection}{Question
1}\label{V-answer-loigeom1}
\(\mathbb{P}(X=n) = p (1-p)^{n-1} = (1-q)q^{n-1} \text{, où } p \text{ est la probabilité de succès, } n=1,2,\ldots\)
\begin{align*}
\mathbb{E}[X]&=\sum_{i=1}^{\infty}i\mathbb{P}(X=i)\\
&=\sum_{i=1}^{\infty} i (1-q)q^{i-1}\\
&=(1-q)\sum_{i=1}^{\infty} iq^{i-1}\\
&=(1-q)\sum_{i=1}^{\infty} \frac{\mathrm{d}}{\mathrm{d} q} q^i\\
&=(1-q)\frac{\mathrm{d}}{\mathrm{d} q}\left(\sum_{i=1}^{\infty}  q^i\right)\\
&=(1-q)\frac{\mathrm{d}}{\mathrm{d} q}\left(\frac{q}{1-q}\right), \text{ car } \sum_{i=1}^{\infty}  q^i = q \sum_{i=0}^{\infty}  q^i = \frac{q}{1-q}\\
&=(1-q)\frac{1}{(1-q)^2} \\
&=\frac{1}{(1-q)}\\
&=\frac{1}{p}
\end{align*} \bigskip Pour la variance
\(\mathbb{V}[X] = \mathbb{E}[X^2]-\mathbb{E}[X]\) \begin{align*}
\mathbb{E}[X^2]&=\sum_{i=1}^{\infty}i^2\mathbb{P}(X=i)\\
&=\sum_{i=1}^{\infty} i^2 (1-q)q^{i-1}\\
&=(1-q)\sum_{i=1}^{\infty} i^2q^{i-1}\\
&=(1-q)\sum_{i=1}^{\infty}\left( i\left((i+1)q^{i-1} - q^{i-1}\right)\right)\\
&=(1-q)\sum_{i=1}^{\infty}\left( i(i+1)q^{i-1} - iq^{i-1}\right)\\
&=(1-q)\sum_{i=1}^{\infty}\left( \frac{\mathrm{d}^2}{\mathrm{d} q^2} q^{i+1} - \frac{\mathrm{d}}{\mathrm{d} q} q^i\right)\\
&=(1-q)\left( \frac{\mathrm{d}^2}{\mathrm{d} q^2}\left(\sum_{i=1}^{\infty}q^{i+1}\right)-\frac{\mathrm{d}}{\mathrm{d} q}\left(\sum_{i=1}^{\infty}  q^i\right)\right)
\end{align*} On a déjà calculé le second terme ci-dessus. \begin{align*}
\frac{\mathrm{d}^2}{\mathrm{d} q^2}\sum_{i=1}^{\infty}q^{i+1} &= \frac{\mathrm{d}^2}{\mathrm{d} q^2}q \sum_{i=1}^{\infty}q^{i}\\
&=\frac{\mathrm{d}^2}{\mathrm{d} q^2}\frac{q^2}{1-q}
\end{align*} Calcul des dérivées~: \begin{align*}
\frac{\mathrm{d}}{\mathrm{d} q}\frac{q^2}{1-q} &= \frac{2q(1-q)+q^2}{(1-q)^2}\\
&= \frac{2q-q^2}{(1-q)^2}
\end{align*}

\begin{align*}
\frac{\mathrm{d}}{\mathrm{d} q}\frac{2q-q^2}{(1-q)^2}&=\frac{2(1-q)(1-q)^2+2(1-q)(2q-q^2)}{(1-q)^4}\\
&= \frac{2(1-q)^2+2(2q-q^2)}{(1-q)^3}\\
&= \frac{2 - 4q +2q^2 + 4q -2q^2}{(1-q)^3}\\
&= \frac{2}{(1-q)^3}
\end{align*} On en déduit \begin{align*}
\mathbb{E}[X^2]&=(1-q)\left(\frac{2}{(1-q)^3}-\frac{1}{(1-q)^2}\right)\\
&=\frac{2}{(1-q)^2}-\frac{1}{(1-q)}\\
&=\frac{2-(1-q)}{(1-q)^2}\\
&=\frac{1+q}{(1-q)^2}
\end{align*} Et finalement \begin{align*}
\mathbb{V}[X] &= \mathbb{E}[X^2]-\mathbb{E}[X]\\
&= \frac{1+q}{(1-q)^2}-\frac{1}{(1-q)^2}\\
&=\frac{q}{(1-q)^2}\\
&=\frac{1-p}{p^2}
\end{align*}

\runinhead{Question
2}\phantomsection\addcontentsline{toc}{subsubsection}{Question
2}\label{V-answer-loigeom2}
cf.~notebook

\runinhead{Question
3}\phantomsection\addcontentsline{toc}{subsubsection}{Question
3}\label{V-answer-loigeom3}

\begin{align*}
\mathbb{P}(\lceil X \rceil = n) & = \mathbb{P}(n-1 < X \leq n) \\
& = F_X(n) - F_X(n-1) \\
& = e^{-\lambda (n-1)}(1-e^{-\lambda})
\end{align*} d'où \(\lambda = -\ln(1-p)\)

\runinhead{Question
4}\phantomsection\addcontentsline{toc}{subsubsection}{Question
4}\label{V-answer-loigeom4}
cf.~notebook

\section{Simulation de la loi gaussienne par la méthode de
rejet}\label{V-simulation-de-la-loi-gaussienne-par-la-muxe9thode-de-rejet-1}

\(X\) est une variable aléatoire de loi gaussienne d'espérance \(m=0\)
et de variance \(\sigma^2=1\), \(\mathcal{N}(0,1)\). Sa densité est
\[f(x) = \frac{1}{\sqrt{2\pi}}\exp\left(-\frac{1}{2}x^2\right),~x\in \mathbb{R}\]

\runinhead{Question
1}\phantomsection\addcontentsline{toc}{subsubsection}{Question
1}\label{V-answer-lgrej1}
Soit \(x \in \mathbb{R}\) \begin{align*}
x^2 & \geq 2|x| - 1 \\
\frac{1}{\sqrt{2\pi}}e^{-x^2/2} & \leq \sqrt{\frac{e}{2\pi}} e^{-|x|}
\end{align*}

\runinhead{Question
2}\phantomsection\addcontentsline{toc}{subsubsection}{Question
2}\label{V-answer-lgrej2}
\(g\) est bien positive et son intégrale sur \(\mathbb{R}\) vaut
\[\int_\mathbb{R}g(x) dx = \int_\mathbb{R}\frac{1}{2} e^{-|x|} dx = \int_{\mathbb{R}^+} e^{-x} dx = 1\]

\runinhead{Question
3}\phantomsection\addcontentsline{toc}{subsubsection}{Question
3}\label{V-answer-lgrej3}
voir dans \protect\hyperref[V-exemples1]{exemples d'application
1 (p.~\pageref*{V-exemples1})}.

\runinhead{Question
4}\phantomsection\addcontentsline{toc}{subsubsection}{Question
4}\label{V-answer-lgrej4}
cf.~notebook

\section{Simulation de la loi de Wigner}\label{V-answer-wigner}

cf.~notebook

\section{Loi des grands nombres et théorème central
limite}\label{V-loi-des-grands-nombres-et-thuxe9oruxe8me-central-limite-1}

\runinhead{Question
1}\phantomsection\addcontentsline{toc}{subsubsection}{Question
1}\label{V-answer-lgntcl1}
cf.~notebook

\runinhead{Question
2}\phantomsection\addcontentsline{toc}{subsubsection}{Question
2}\label{V-answer-lgntcl2}
cf.~notebook

\runinhead{Question
3}\phantomsection\addcontentsline{toc}{subsubsection}{Question
3}\label{V-answer-lgntcl3}
cf.~notebook

\section{Simulation d'un mélange de
gaussiennes}\label{V-simulation-dun-muxe9lange-de-gaussiennes-1}

\runinhead{Question
1}\phantomsection\addcontentsline{toc}{subsubsection}{Question
1}\label{V-answer-melloi1}
cf.~notebook

\runinhead{Question
2}\phantomsection\addcontentsline{toc}{subsubsection}{Question
2}\label{V-answer-melloi2}
cf.~notebook

\section{Echantillonnage
d'importance}\label{V-echantillonnage-dimportance-3}

\runinhead{Question
1}\phantomsection\addcontentsline{toc}{subsubsection}{Question
1}\label{V-answer-echimp1}
On a
\(F_{X|X>3}(x) = \mathbb{P}(X \leq x |X > 3) = \frac{\mathbb{P}(X\leq x, X>3)}{1-F_X(3)} = \frac{F_X(x)-F_X(3)}{1-F_X(3)}1_{]3,+\infty[}(x)\)
d'où \(f_{X|X>3}(x) = \frac{f_X(x)}{1-F_X(3)}1_{]3,+\infty[}(x)\)

\runinhead{Question
2}\phantomsection\addcontentsline{toc}{subsubsection}{Question
2}\label{V-answer-echimp2}
cf.~notebook

\runinhead{Question
3}\phantomsection\addcontentsline{toc}{subsubsection}{Question
3}\label{V-answer-echimp3}
cf.~notebook
